% Part: many-valued-logic
% Chapter: syntax-and-semantics
% Section: connectives

\documentclass[../../../include/open-logic-section]{subfiles}

\begin{document}

\olfileid[id]{mvl}{syn}{con}

\olsection{Bahasa dan Konektif}

Logika proposisional klasik, dan banyak logika lainnya, menggunakan suatu
persediaan tetap \emph{konstanta proposisional} dan \emph{konektif}.
Sebagai contoh, kita menggunakan unsur-unsur berikut sebagai unsur primitif:
\begin{enumerate}
\tagitem{prvFalse}{Konstanta proposisional untuk !!{falsity}~$\lfalse$.}{}
\tagitem{prvTrue}{Konstanta proposisional untuk !!{truth}~$\ltrue$.}{}
\item Konektif-konektif logis:
  \startycommalist
  \iftag{prvNot}{\ycomma $\lnot$ (negasi)}{}%
  \iftag{prvAnd}{\ycomma $\land$ (konjungsi)}{}%
  \iftag{prvOr}{\ycomma $\lor$ (disjungsi)}{}%
  \iftag{prvIf}{\ycomma $\lif$ (!!{conditional})}{}%
  \iftag{prvIff}{\ycomma $\liff$ (!!{biconditional})}{}%
\end{enumerate}
\iftag{defNot,defOr,defAnd,defIf,defIff,defTrue,defFalse,defEx,defAll}{%
Selain konektif primitif di atas, kita juga menggunakan simbol yang
didefinisikan sebagai singkatan, seperti
\startycommalist
  \iftag{defNot}{\ycomma $\lnot$ (negasi)}{}%
  \iftag{defAnd}{\ycomma $\land$ (konjungsi)}{}%
  \iftag{defOr}{\ycomma $\lor$ (disjungsi)}{}%
  \iftag{defIf}{\ycomma $\lif$ (!!{conditional})}{}%
  \iftag{defIff}{\ycomma $\liff$ (!!{biconditional})}{}%
  \iftag{defFalse}{\ycomma $\lfalse$ (!!{falsity})}{}%
  \iftag{defTrue}{\ycomma $\ltrue$ (!!{truth})}.}{}

Konektif yang sama juga digunakan dalam logika bernilai banyak. Akan tetapi,
sering kali berguna untuk menyertakan beberapa versi, misalnya, konjungsi
dalam logika yang sama; hal ini memerlukan simbol-simbol yang berbeda agar
versi-versi tersebut tetap terpisah. Beberapa logika bernilai banyak juga
menyertakan konektif yang tidak memiliki padanan dalam logika klasik. Karena
itu, kita akan menggunakan kerangka yang sedikit lebih umum daripada biasanya.

\begin{defn}
  Suatu \emph{bahasa proposisional} terdiri atas suatu himpunan $\Lang L$
  yang beranggotakan \emph{konektif}. Setiap konektif $\star$ mempunyai suatu
  \emph{aritas}; konektif beraritas~$n$ disebut konektif \emph{$n$-tempat}.
  Konektif beraritas~$0$ juga disebut \emph{konstanta}; konektif
  beraritas~$1$ disebut \emph{uner}, sedangkan konektif beraritas~$2$
  disebut \emph{biner}.
\end{defn}

\begin{ex}
  Bahasa baku logika proposisional~$\Lang L_0$ terdiri atas konektif-konektif
  berikut (beserta aritasnya):
  $\lfalse$~($0$),
  $\lnot$~($1$),
  $\land$~($2$),
  $\lor$~($2$),
  $\lif$~($2$). Kebanyakan logika yang kita bahas menggunakan bahasa ini.
  Menurut tradisi dan konvensi, beberapa logika menggunakan simbol yang
  berbeda bagi konektif tertentu. Sebagai contoh, dalam logika produk,
  konjungsi sering dilambangkan dengan $\odot$, bukan~$\land$. Kadang-kadang
  berguna pula untuk menambahkan operator baru, misalnya operator ketertentuan
  $\triangle$ ($1$-tempat).
\end{ex}

\end{document}
